\documentclass[11pt,a4paper]{article}

% ---------------------------------------------------------------
% Packages
% ---------------------------------------------------------------
\usepackage[utf8]{inputenc}
\usepackage[T1]{fontenc}
\usepackage{lmodern}
\usepackage[margin=2.5cm]{geometry}
\usepackage{amsmath,amssymb}
\usepackage{booktabs}
\usepackage{tabularx}
\usepackage{array}
\usepackage{enumitem}
\usepackage[numbers,sort&compress]{natbib}
\usepackage[hidelinks]{hyperref}

\newcolumntype{L}{>{\raggedright\arraybackslash}X}

\title{Reproducing and Assessing OTPWR:\\
Globally Optimal, Dynamically Feasible Trajectory Planning\\
for Wheeled Mobile Robots}
\author{}
\date{\today}

\begin{document}
\maketitle

% ===============================================================
\section{Introduction}
\label{sec:intro}
% ===============================================================

Demand for round-the-clock monitoring of buildings, industrial campuses,
substations and chemical plants keeps growing, while the security and
inspection workforce is ageing and labour costs are rising
\citep{luo2024otpwr,gao2018patrol}. Wheeled mobile robots that can patrol
such sites on their own are an attractive answer, and their usefulness
depends heavily on \emph{motion planning}: once a map of the environment is
available, the robot must compute a feasible trajectory from its current pose
to a target pose that is collision-free, respects the robot's kinematic and
dynamic limits, and is good with respect to some criterion such as path
length, travel time or energy
\citep{lavalle2006planning,siegwart2011introduction}.

In practice, motion planning for ground robots is almost always split into
two layers \citep{marder2010office,paden2016survey}. A \emph{global planner}
searches a map for a sequence of collision-free waypoints. Typical choices
are graph search on an occupancy grid (Dijkstra, A$^\ast$) or sampling-based
planners (PRM, RRT$^\ast$). A \emph{local planner}, such as the Dynamic Window
Approach (DWA) \citep{fox1997dwa}, artificial potential fields
\citep{khatib1986realtime} or the Timed Elastic Band (TEB)
\citep{rosmann2012teb}, then turns those waypoints into velocity commands
over a short horizon while respecting the robot's dynamics. This split makes
both layers tractable, but it has a known weakness. The global planner
optimises a purely geometric cost and ignores velocity and acceleration, so
its output is a list of points rather than a trajectory. The local planner
does account for dynamics, but it optimises a short-horizon objective and can
therefore get stuck in local minima, oscillate, or deviate from (or even skip)
global waypoints. The motion that results is neither globally optimal nor
guaranteed to be smooth.

Luo et al.\ \citep{luo2024otpwr} propose \emph{Optimal Trajectory Planning for
Wheeled Robots} (OTPWR) to close this gap with a single pipeline that has
three stages:
\begin{enumerate}[label=(\roman*),nosep]
  \item \textbf{Dimensionality reduction by differential flatness.} The
        unicycle model of a differential-drive robot is differentially flat
        with flat output $\sigma=[x,y]^{\mathsf T}$
        \citep{fliess1995flatness,murray1995flat}. The heading $\theta$ and
        the inputs $(v,\omega)$ are recovered algebraically from
        $\dot x,\dot y,\ddot x,\ddot y$, so planning can take place in the
        two-dimensional plane.
  \item \textbf{Global search with a diagonal-distance A$^\ast$.} An
        eight-connected occupancy grid is searched with A$^\ast$ guided by
        the octile (diagonal) heuristic and a small tie-breaking factor
        $(1+p)$. This yields a shortest collision-free waypoint sequence.
  \item \textbf{Minimum-snap trajectory optimisation.} Each flat output is
        represented as a piecewise fifth-order polynomial through the
        waypoints. The integral of the squared fourth derivative (snap) is
        minimised as a convex quadratic program (QP), subject to waypoint
        interpolation, velocity/acceleration continuity, and bounds on
        velocity and acceleration, and is solved with OSQP
        \citep{stellato2020osqp}.
\end{enumerate}
In simulation the authors report that diagonal A$^\ast$ finds the same
shortest path as Dijkstra's algorithm while cutting search time by 46.6\%
and the number of visited nodes by 43.1\%. They also report that the
resulting trajectory is smoother than a DWA baseline and passes through every
waypoint while DWA misses one. Experiments on a hub-motor robot running ROS
complete the evaluation.

These claims are attractive, but several details that a reproduction needs
are not stated in the article. Among them are the segment-time allocation
$T_j$, the tie-breaking factor $p$, the step costs $D_{\mathrm{straight}}$ and
$D_{\mathrm{diagonal}}$, and the exact benchmark maps. A second issue is
structural. The QP enforces velocity and acceleration bounds only at the
segment junctions, and it includes no constraint that keeps the polynomial
inside free space \emph{between} waypoints. ``Globally optimal'' therefore
holds only for the fixed waypoints and segment times. It does not hold for
the continuous-time trajectory as a whole.

This report reproduces OTPWR independently and examines it critically. We
ask four questions:
\begin{enumerate}[label=\textbf{Q\arabic*.},nosep]
  \item Can the published path-search results (Tables~1--2 of
        \citep{luo2024otpwr}) be reproduced?
  \item Do the conclusions hold across many randomly generated grid maps and
        across maps built with SLAM in simulation?
  \item Does the trajectory respect the velocity and acceleration limits and
        stay collision-free when checked densely in continuous time, and not
        only at the waypoints?
  \item Does OTPWR outperform DWA on the same waypoints in smoothness and
        waypoint attainment?
\end{enumerate}

The contributions of this work are:
\begin{itemize}[nosep]
  \item a Python implementation of the full OTPWR pipeline, together
        with Dijkstra, Euclidean, Manhattan and diagonal A$^\ast$ baselines
        and a DWA baseline ported from a reference MATLAB implementation;
  \item a documented list of the assumptions needed to fill the gaps in the
        article, each tested at several values;
  \item a continuous-time audit of dynamic feasibility and obstacle clearance
        along the optimised trajectories;
  \item a discussion of extensions, namely convex safe corridors,
        whole-segment bound enforcement and time-allocation optimisation,
        that would make the method's optimality and safety claims hold in
        continuous time.
\end{itemize}

The rest of the report is organised as follows.
Section~\ref{sec:lit} reviews the literature on global search,
sampling-based and local planning, differential flatness, and
optimisation-based trajectory generation, and places OTPWR within it.
Later sections present the methodology, the experimental results and the
discussion.

% ===============================================================
\section{Literature Review}
\label{sec:lit}
% ===============================================================

Over the past decades, motion planning for wheeled robots has moved from
coarse path planning over isolated points in space towards fine trajectory
planning in time, helped by growing on-board computing power
\citep{luo2024otpwr,gonzalez2016review}. This review follows that
progression. It covers graph-based global search
(Section~\ref{sec:lit-graph}), sampling-based and kinodynamic planning
(Section~\ref{sec:lit-sampling}), local and reactive planners
(Section~\ref{sec:lit-local}), differential flatness
(Section~\ref{sec:lit-flat}), optimisation-based polynomial trajectory
generation (Section~\ref{sec:lit-poly}), and continuous-time safety and
feasibility (Section~\ref{sec:lit-safe}). Section~\ref{sec:lit-gap}
summarises the research gap.

% ---------------------------------------------------------------
\subsection{Graph-based global path search}
\label{sec:lit-graph}

Graph search needs a discrete representation of the workspace. Common choices
are visibility graphs, Voronoi diagrams and occupancy grids
\citep{lavalle2006planning}. Visibility graphs give shortest paths that touch
obstacle corners. Voronoi diagrams maximise clearance at the cost of path
length. Grid maps divide the plane at a fixed resolution and connect cells in
four or eight directions. They are by far the most common choice for indoor
robots because they map directly onto the occupancy grids produced by 2-D
LiDAR SLAM.

Dijkstra's algorithm \citep{dijkstra1959note} expands nodes in order of their
cost-to-come $g(n)$ and returns an optimal path on any graph with
non-negative edge weights. Because it has no sense of where the goal lies, it
expands nodes uniformly in all directions. A$^\ast$ \citep{hart1968formal}
orders expansion by $f(n)=g(n)+h(n)$. If the heuristic $h$ is
\emph{admissible} ($h(n)\le h^\ast(n)$), A$^\ast$ still returns an optimal
path, and if $h$ is also consistent, no node is expanded more than once. The
heuristic must suit the connectivity of the grid. On an eight-connected grid,
the Manhattan distance overestimates the true cost and so loses optimality.
The Euclidean distance is admissible but loose. The octile, or diagonal,
distance
\begin{equation}
  h(n)=D_{\mathrm{s}}\,(d_x+d_y)+(D_{\mathrm{d}}-2D_{\mathrm{s}})\min(d_x,d_y)
\end{equation}
equals the exact obstacle-free cost and is therefore the tightest admissible
choice. This is the heuristic OTPWR adopts. Many nodes on such grids have
equal $f$ values, so tie-breaking strongly affects how many nodes are
expanded. Multiplying $h$ by a small factor $(1+p)$ favours nodes nearer the
goal, as OTPWR does. In doing so it moves towards weighted A$^\ast$
\citep{pohl1970heuristic}, which trades a bounded loss of optimality for
speed. Anytime variants such as ARA$^\ast$ \citep{likhachev2003ara} make that
trade-off adjustable at run time.

Several extensions target speed or changing environments. Jump Point Search
(JPS) \citep{harabor2011online} prunes symmetric paths on uniform-cost grids
and expands only ``jump points'', which can speed search up by an order of
magnitude while keeping optimality. Fan and Guo \citep{fan2023improved}
adapt JPS for a seabed mining vehicle. D$^\ast$ \citep{stentz1994optimal} and
D$^\ast$~Lite \citep{koenig2002dstar} repair the solution incrementally when
the map changes, which matters for robots that discover obstacles as they
move. Recent application papers keep refining A$^\ast$ for particular
domains, for example orchard mowers \citep{zhang2024orchard} and multimodal
urban transport \citep{feng2023improved}. Liu et al.\ \citep{liu2023kinematically}
add kinematic constraints to A$^\ast$ and fuse it with DWA.

All of these methods share one limitation for our purposes: they return a
sequence of grid cells. The path has corners where the heading changes
abruptly, and it carries no information about velocity or time, so it cannot
be executed directly by a robot with bounded acceleration.

% ---------------------------------------------------------------
\subsection{Sampling-based and kinodynamic planning}
\label{sec:lit-sampling}

Sampling-based planners avoid discretising the space explicitly and scale to
high-dimensional configuration spaces. Probabilistic Roadmaps (PRM)
\citep{kavraki1996probabilistic} build a reusable graph of random collision-free
samples. Rapidly-exploring Random Trees (RRT) \citep{lavalle1998rapidly}
grow a tree from the start and are probabilistically complete.
Karaman and Frazzoli \citep{karaman2011sampling} showed that RRT almost
surely converges to a suboptimal solution and introduced RRT$^\ast$ and
PRM$^\ast$, which are asymptotically optimal. Informed RRT$^\ast$
\citep{gammell2014informed} restricts sampling to an ellipsoid once a first
solution exists and so speeds up convergence. Recent variants include
double-layer, goal-biased anytime RRT$^\ast$ \citep{esmaiel2024double}, an
improved RRT$^\ast$ for disinfection robots \citep{wang2024improved}, and
Informed-Bi-Quick RRT$^\ast$ for dual-arm systems \citep{zhang2024informed}.

Plain RRT$^\ast$ joins nodes with straight lines and so ignores the system's
dynamics. \emph{Kinodynamic} planners join nodes with dynamically feasible
local trajectories instead \citep{lavalle2001randomized}. LQR-RRT$^\ast$
\citep{perez2012lqr} uses a linear-quadratic regulator to derive the
extension heuristic and distance metric. Kinodynamic RRT$^\ast$
\citep{webb2013kinodynamic} solves a fixed-final-state, free-final-time
optimal control problem between nodes of a linear system. For differentially
flat systems, the steering problem can instead be written as an optimisation
over the flat outputs (Section~\ref{sec:lit-flat}), which yields optimal
motion primitives between states. On the search-based side, state lattices
\citep{pivtoraiko2009differentially} and Hybrid A$^\ast$
\citep{dolgov2010path} expand nodes with kinematically feasible motion
primitives, so the path they return can be driven.

These methods couple search and dynamics, but at a price. Sampling-based
planners give only probabilistic guarantees and converge slowly in practice.
Kinodynamic steering must be solved at every edge, which is expensive. For a
planar robot on a 2-D grid, where the search space is small, systematic graph
search remains the more efficient choice. OTPWR follows this reasoning: it
keeps a cheap 2-D graph search and adds dynamics afterwards in a separate
optimisation stage.

% ---------------------------------------------------------------
\subsection{Local and reactive planners}
\label{sec:lit-local}

Local planners produce the robot's commands over a short horizon. The
Artificial Potential Field (APF) method \citep{khatib1986realtime} treats the
goal as an attractive field and obstacles as repulsive fields. It is simple
and fast, but it has well-known local minima and oscillates in narrow
passages \citep{koren1991potential}. Variants such as APF-S2T
\citep{chen2024apf} are still being developed. The Dynamic Window Approach
\citep{fox1997dwa} samples linear and angular velocities $(v,\omega)$ within
the window reachable under the acceleration limits, simulates each pair
forward, and picks the one that maximises a weighted sum of heading,
clearance and speed. DWA respects the robot's dynamics by construction and
is the default local planner in the ROS navigation stack
\citep{marder2010office}. Recent work adds energy management and fuzzy-logic
weight tuning \citep{yao2024improved}. Elastic bands
\citep{quinlan1993elastic} deform a global path as if it were under internal
and external forces. The Timed Elastic Band (TEB)
\citep{rosmann2012teb,rosmann2017integrated} adds time intervals between
poses and solves a sparse nonlinear least-squares problem over kinematic,
dynamic and clearance objectives. It has been applied to cooperative
steering robots \citep{sun2022solution}.

These local methods share the weaknesses that motivate OTPWR. Their
objectives are local, so the motion that results is not globally optimal.
Sampling-based methods such as DWA produce piecewise-constant, jagged
velocity profiles. When the local goal sequence is followed with a finite
look-ahead window, a waypoint can be overshot or skipped. Fusing A$^\ast$
with DWA \citep{liu2023kinematically} keeps these layers separate rather than
unifying them.

% ---------------------------------------------------------------
\subsection{Differential flatness}
\label{sec:lit-flat}

A nonlinear system is \emph{differentially flat} if there is a set of outputs,
equal in number to the inputs, such that every state and input can be written
as an algebraic function of these outputs and finitely many of their
derivatives \citep{fliess1995flatness}. Flatness turns trajectory generation
into a problem in the space of flat outputs. Any sufficiently smooth flat
output curve corresponds to a feasible state--input trajectory, so the
system's differential constraints no longer need to be enforced explicitly.
Murray et al.\ \citep{murray1995flat} catalogue flat mechanical systems,
including the unicycle and the kinematic car. Ryu and Agrawal
\citep{ryu2011differential} use flatness for robust control of mobile robots
that slip.

For the unicycle model
\begin{equation}
  \dot x = v\cos\theta,\qquad
  \dot y = v\sin\theta,\qquad
  \dot\theta = \omega,
\end{equation}
the position $\sigma=[x,y]^{\mathsf T}$ is a flat output, with
\begin{equation}
  \theta=\operatorname{atan2}(\dot y,\dot x),\qquad
  v=\sqrt{\dot x^2+\dot y^2},\qquad
  \omega=\frac{\dot x\ddot y-\dot y\ddot x}{\dot x^2+\dot y^2}.
\end{equation}
The map is singular when $v=0$, so the heading at rest, and at the start and
goal in particular, must be handled separately. Flatness is also the basis
of minimum-snap planning for quadrotors \citep{mellinger2011minimum}, where
the flat outputs are position and yaw. OTPWR applies the same idea to a
planar wheeled robot.

% ---------------------------------------------------------------
\subsection{Optimisation-based polynomial trajectory generation}
\label{sec:lit-poly}

Once planning happens in flat-output space, a natural trajectory model is a
piecewise polynomial through the waypoints. The idea of minimising the
integral of a squared derivative goes back to the minimum-jerk model of human
arm motion \citep{flash1985coordination}. Mellinger and Kumar
\citep{mellinger2011minimum} formulated minimum-snap trajectory generation
for quadrotors as a QP over polynomial coefficients. The cost
$\int\!\big(f^{(k)}(t)\big)^2\,dt$ is a convex quadratic form
$\mathbf p^{\mathsf T}H\mathbf p$ in the coefficients $\mathbf p$, and
waypoint, continuity and boundary conditions are linear equalities. Richter,
Bry and Roy \citep{richter2016polynomial} showed that the unconstrained
problem can be solved in closed form by optimising over endpoint derivatives,
which improves numerical conditioning. They also used an RRT$^\ast$ path as
the waypoint sequence and optimised the segment times by gradient descent.
Oleynikova et al.\ \citep{oleynikova2016continuous} added collision costs for
online replanning, and Preiss et al.\ \citep{preiss2017downwash} extended
polynomial planning to large quadrotor teams. For road vehicles, the Baidu
Apollo EM planner \citep{fan2018baidu} splits planning into path and speed
sub-problems, each solved by dynamic programming followed by spline QP.

More general trajectory optimisation methods such as CHOMP
\citep{ratliff2009chomp} and TrajOpt \citep{schulman2014motion} optimise over
time-discretised states with collision terms. They handle complex geometry
but are non-convex and depend on initialisation. The polynomial QP is convex,
because $H$ is a Gram matrix and is positive semidefinite
\citep{boyd2004convex}, so it has a global optimum for fixed waypoints and
segment times. It can be solved in milliseconds by modern first-order solvers
such as OSQP \citep{stellato2020osqp}. OTPWR adopts this formulation with
$k=4$ (snap) and fifth-order segments. It adds linear inequality constraints
on velocity and acceleration at the segment junctions.

% ---------------------------------------------------------------
\subsection{Continuous-time safety, feasibility and time allocation}
\label{sec:lit-safe}

Waypoint-constrained polynomial QPs have three known gaps, and the aerial
robotics literature has addressed each of them.

\paragraph{Collision between waypoints.} A polynomial that interpolates
collision-free waypoints can still cut corners or overshoot into obstacles.
Richter et al.\ \citep{richter2016polynomial} insert extra waypoints and
re-solve until the trajectory is collision-free. A more principled approach
is to build a sequence of convex free-space regions (a \emph{safe flight
corridor}) around the discrete path and to constrain each segment to stay
inside its region. Examples are axis-aligned cubes \citep{chen2016online},
convex polyhedra obtained by ellipsoid inflation \citep{liu2017planning}, and
convex regions found by IRIS together with mixed-integer segment assignment
\citep{deits2015efficient}.

\paragraph{Bounds inside segments.} Imposing $|v|\le v_{\max}$ and
$|a|\le a_{\max}$ only at the junctions does not bound the extrema inside a
segment. Gao et al.\ \citep{gao2018online} write each segment in the
Bernstein basis. Because of the convex-hull property, linear constraints on
the control points then guarantee that position, velocity and acceleration
stay within bounds over the whole segment. Zhou et al.\
\citep{zhou2019robust} use B-splines with the same property inside a fast
kinodynamic planner.

\paragraph{Time allocation.} With fixed segment durations $T_j$ the QP is
convex, but the overall optimum depends strongly on the $T_j$. Heuristics
such as dividing distance by maximum speed can be far from optimal.
Gradient-based refinement \citep{richter2016polynomial} and joint
spatial--temporal optimisation, such as MINCO \citep{wang2022geometrically},
optimise the durations directly while keeping the inner problem efficient.

% ---------------------------------------------------------------
\subsection{Summary and research gap}
\label{sec:lit-gap}

Table~\ref{tab:lit} places the main approaches side by side. Graph search
gives global optimality cheaply in 2-D, but its output is a list of points.
Local planners respect dynamics only locally. Optimisation-based polynomial
methods give smooth, dynamically consistent trajectories, but they need good
waypoints, good time allocation and explicit safety constraints to be
optimal and safe in continuous time.

\begin{table}[t]
  \centering
  \caption{Comparison of motion-planning families relevant to OTPWR.}
  \label{tab:lit}
  \small
  \begin{tabularx}{\textwidth}{@{}l L L L@{}}
    \toprule
    \textbf{Family} & \textbf{Representative work} & \textbf{Strength} & \textbf{Limitation for wheeled robots} \\
    \midrule
    Graph search &
      Dijkstra \citep{dijkstra1959note}, A$^\ast$ \citep{hart1968formal}, JPS \citep{harabor2011online}, D$^\ast$~Lite \citep{koenig2002dstar} &
      Complete and optimal on the grid; fast in 2-D &
      Discrete waypoints; no velocity or time; sharp corners \\
    \addlinespace
    Sampling-based &
      PRM \citep{kavraki1996probabilistic}, RRT$^\ast$ \citep{karaman2011sampling}, Informed RRT$^\ast$ \citep{gammell2014informed} &
      Scales to high dimensions; asymptotically optimal &
      Only probabilistic guarantees; slow convergence; jagged paths \\
    \addlinespace
    Kinodynamic &
      LQR-RRT$^\ast$ \citep{perez2012lqr}, Kinodynamic RRT$^\ast$ \citep{webb2013kinodynamic}, Hybrid A$^\ast$ \citep{dolgov2010path} &
      Dynamically feasible edges &
      Costly steering problem at every edge \\
    \addlinespace
    Local / reactive &
      APF \citep{khatib1986realtime}, DWA \citep{fox1997dwa}, TEB \citep{rosmann2017integrated} &
      Real-time; respects dynamics locally &
      Local minima; not globally optimal; may skip waypoints \\
    \addlinespace
    Polynomial QP &
      Minimum snap \citep{mellinger2011minimum,richter2016polynomial}, Apollo EM \citep{fan2018baidu} &
      Convex; smooth; millisecond solve times &
      Needs waypoints and $T_j$; safety only at waypoints \\
    \addlinespace
    Corridor / Bernstein &
      SFC \citep{liu2017planning}, Bernstein \citep{gao2018online}, MINCO \citep{wang2022geometrically} &
      Continuous-time safety and bounds &
      More constraints; corridor construction cost \\
    \bottomrule
  \end{tabularx}
\end{table}

OTPWR \citep{luo2024otpwr} combines three established ideas: flatness-based
reduction to 2-D, an optimal graph search with a tight heuristic, and a
minimum-snap QP. Together they replace the decoupled global/local
architecture for wheeled robots with a single global trajectory. The idea of
each stage is not new, since minimum-snap planning over a discrete path is
standard practice for aerial robots \citep{richter2016polynomial}. The
contribution of OTPWR is to apply this combination to differential-drive
ground robots and to compare it with DWA. Several questions remain open:
\begin{itemize}[nosep]
  \item \textbf{Reproducibility.} Key parameters ($p$, $D_{\mathrm{s}}$,
        $D_{\mathrm{d}}$, $T_j$, the benchmark maps) are not reported. The
        reported speed-ups come from only two maps and single runs, and
        several in-text citations in the article do not match its reference
        list.
  \item \textbf{Continuous-time feasibility.} Velocity and acceleration are
        bounded only at the junctions $T_j$, and no corridor constraint keeps
        the segments in free space. The safety-corridor and Bernstein-basis
        techniques reviewed above are not used.
  \item \textbf{Scope of optimality.} The trajectory is globally optimal
        only with respect to the fixed waypoints and fixed $T_j$. The
        waypoints are optimal for a geometric cost, not a dynamic one.
  \item \textbf{Generality.} The evaluation does not include randomised map
        sets or SLAM-built maps.
\end{itemize}
This project addresses these questions. It reproduces the published
benchmarks, tests the method across many grid and SLAM maps, audits the
trajectories densely in continuous time, and outlines how corridor
constraints, whole-segment bounds and time-allocation optimisation could make
OTPWR's claims hold in continuous time.

% ---------------------------------------------------------------
\bibliographystyle{unsrtnat}
\bibliography{references}

\end{document}
